\documentclass[11pt]{article}
\pdfinfoomitdate=1
\pdftrailerid{}
\pdfsuppressptexinfo=15
\input{glyphtounicode-cmex}
\pdfgentounicode=1
\usepackage[T1]{fontenc}
\usepackage{lmodern}
\usepackage[margin=1in]{geometry}
\usepackage{amsmath,amssymb,amsthm,mathtools}
\usepackage{microtype}
\usepackage{needspace}
\usepackage[colorlinks=true,linkcolor=blue,citecolor=blue,urlcolor=blue]{hyperref}
\usepackage[nameinlink,noabbrev]{cleveref}
\hypersetup{pdfauthor={OpenAI},pdftitle={Ergodicity of triangular billiards with an irrational angle}}
\newtheorem{theorem}{Theorem}
\newtheorem{corollary}[theorem]{Corollary}
\newtheorem{lemma}[theorem]{Lemma}
\newtheorem{proposition}[theorem]{Proposition}
\theoremstyle{remark}
\newtheorem{remark}[theorem]{Remark}
\newcommand{\R}{\mathbb R}
\newcommand{\C}{\mathbb C}
\newcommand{\Z}{\mathbb Z}
\newcommand{\Q}{\mathbb Q}
\newcommand{\dd}{\,d}
\newcommand{\ind}{\mathbf 1}
\newcommand{\eps}{\varepsilon}
\newcommand{\pa}{\partial}
\newcommand{\pb}{\bar\partial}
\newcommand{\norm}[1]{\left\lVert#1\right\rVert}
\newcommand{\ip}[2]{\left\langle#1,#2\right\rangle}
\setlength{\emergencystretch}{2em}
\title{Ergodicity of triangular billiards\\with an irrational angle}
\author{OpenAI}
\date{September 25, 2026}

\begin{document}
\maketitle

\begin{abstract}
We prove that the unit-speed billiard flow in every nondegenerate
Euclidean triangle with at least one angle irrational relative to
\(\pi\) is ergodic for normalized area times uniform angular measure.
No genericity or Diophantine condition is required. The flow is considered
outside the null set of trajectories that hit a vertex.
\end{abstract}

\section{Introduction}

A billiard trajectory in a Euclidean triangle moves at unit speed and
reflects specularly at an open side: its tangential velocity is preserved
and its normal velocity changes sign. The natural probability measure
is normalized area in the table times uniform angular measure.
Ergodicity asks whether every measurable set invariant under this flow
has probability zero or one. Straight sides supply no dispersing
curvature, while trajectories meeting a vertex have no prescribed
continuation. The effect of repeated reflections therefore carries the
essential dynamical difficulty.

Straightening reflected trajectories by unfolding connects polygonal
billiards to flat geometry; the phase-space construction and null
vertex-hit exception appear explicitly in Zemlyakov--Katok \cite{ZK}.
For a rational polygon, whose angles are rational multiples of \(\pi\),
the finite reflection group separates the full phase space into invariant
directional surfaces. Kerckhoff, Masur, and Smillie proved unique
ergodicity in almost every direction \cite{KMS}; see also the account of
Masur--Tabachnikov \cite[Sections~3.2 and~3.5]{MT}.
The same work yields full-flow ergodicity for a dense \(G_\delta\)
family of tables. Its rational-approximation and Baire-category argument
does not assert ergodicity for each irrational triangle.

Vorobets gave a quantitative sufficient approximation condition and
explicit ergodic examples \cite[Definition~1.3 and Theorem~1.1]{Vorobets}.
The condition requires exceptionally fast simultaneous rational
approximation of the angles; his examples include irrational right
triangles. More recently, Chaika and Forni proved weak mixing of the full
billiard flow for a dense \(G_\delta\) set of \(n\)-gons for each fixed
\(n\ge3\) \cite{CF}. These results distinguish generic
or specially approximable tables from a conclusion for every triangle
having an irrational angle. Numerical studies of irrational right
triangles and fully irrational isosceles triangles have raised doubts
about ergodicity \cite{WCP,ZSBJ}. Their finite-time evidence does not
establish a positive-measure invariant set and leaves room for very
slow convergence of averages.

\Needspace{18\baselineskip}
\begin{theorem}\label{thm:main}
Let \(Q\) be the interior of a nondegenerate Euclidean triangle with angles
\(\alpha,\beta,\gamma\). Suppose that at least one of
\[
 \alpha/\pi,\qquad \beta/\pi,\qquad \gamma/\pi
\]
is irrational.
The unit-speed billiard flow \(\Phi_t\), with specular reflection at the
open sides, is ergodic on \(Q\times S^1\) for the probability measure
\[
 d\mu=\frac{dA\,d\theta}{2\pi\operatorname{Area}(Q)}.
\]
Here trajectories hitting a vertex in either time direction are discarded.
In particular, if a measurable set \(A\) satisfies
\(\mu(\Phi_t^{-1}A\mathbin\triangle A)=0\) for every \(t\in\R\), then
\(\mu(A)\in\{0,1\}\).
\end{theorem}

In particular, the theorem applies to every triangle for which all three
angle ratios are irrational. It also includes triangles with one rational
angle and two irrational angles. No genericity or Diophantine condition
is imposed. The conclusion concerns the full unit-speed phase space;
it is an ergodicity statement for its specified probability measure,
not unique ergodicity, mixing, or a claim about each individual orbit.

For each fixed triangle in Theorem~\ref{thm:main}, the classical
conclusion also gives interior and boundary quantum ergodicity for every
Dirichlet or Neumann eigenbasis. For every real such eigenbasis, the
number of nodal domains tends to infinity along a density-one
subsequence. We state the precise spectral conclusions and verify their
hypotheses in \Cref{sec:spectral}.

The main issue is low regularity. Forni and Moll provide an analytic
framework for flat surfaces with cone points \cite{FM}. Their angular
Fourier decomposition and Cauchy--Riemann identities relate transport
invariance to a recurrence among coefficient derivatives. Under
non-rational holonomy, their Theorem~4.1 proves constancy of an
\(L^2\) function invariant under both horizontal directions. Their
Corollary~4.2 proves constancy of a geodesic-invariant function in the
horizontal Sobolev space \(H^{1,0}\), defined by completion of a
cone-adapted smooth domain with square-integrable derivatives in both
horizontal directions \cite[Section~3.4]{FM}. A bounded invariant
indicator has no such regularity a priori.
We work directly with bounded invariant functions, including invariant
indicators. The localized energy estimate in \Cref{sec:energy} first
gives an \(L^2\) gradient for each angular coefficient. This is a bound
on individual coefficients, not a square-summable family of gradients.
The remaining step is to prove that each such gradient vanishes.

\paragraph{Proof strategy.}
Glue two copies of the triangle along corresponding sides to form its
oriented flat double, whose singular points lie at the three vertices.
On its unit tangent bundle, let \(X\) denote
differentiation along a geodesic and \(Y\) differentiation in the
perpendicular parallel direction. The equation \(Xf=0\) relates
neighboring angular Fourier coefficients of an invariant function \(f\).
For each integer \(m\), the sum of angular modes \(m,m+2,m+4,\ldots\)
has an \(X\)-derivative consisting of just one endpoint coefficient
derivative. We must show that this endpoint is zero.

We smooth only in position, away from the three cone points. Averaging
over long segments that stay clear of the cones gives bounded test
functions close to \(f\), with a small transport derivative. Integrating
the smoothed endpoint equation along these segments yields a pairing
that tends to zero. The energy estimate controls the two coefficient
gradients in that pairing while the smoothing scale tends to zero.
The resulting weak limit satisfies the exact local transport equation.
Applying the bounded-coefficient estimate again supplies a new bound
independent of the remaining approximation parameter. Only then do we
let the weak limits approach \(f\). This second passage proves that the
endpoint derivative, and hence every coefficient gradient, vanishes.
No horizontal Sobolev hypothesis is imposed on \(f\).

Parallel coefficients must respect the rotation acquired around a cone
point. A single irrational angle eliminates all nonzero angular modes;
connectedness makes the remaining mode constant. This isolates the
angle hypothesis from the analytic argument.

All the analytic steps needed for this passage are proved below. Constants
may depend on the fixed triangle, but not on Fourier indices or small
smoothing parameters.

\Cref{sec:geometry} sets up the double and coefficient bundles.
\Cref{sec:smoothing} supplies the supported smoothing identities, and
\Cref{sec:energy} proves the individual-coefficient estimates.
\Cref{sec:rigidity} carries out the two successive limits.
\Cref{sec:conclusion} completes the holonomy argument, and
\Cref{sec:spectral} derives the spectral consequences.

\section{The double and angular coefficients}\label{sec:geometry}

We pass to a surface on which reflection becomes straight continuation.
Throughout the analytic argument, the triangle is fixed and
nondegenerate; no irrationality assumption is needed until
\Cref{sec:conclusion}.

Glue two closed copies of the triangle along corresponding sides, give
the copies opposite orientations, and remove the three vertices. The
result is a connected oriented flat surface \(M\). Its completion is
compact and has cone angles \(2\alpha,2\beta,2\gamma\) at the finite set
\(V\) of removed points. Open edges are regular: continuation across
an edge is described by reflecting a Euclidean chart on one copy.
Transitions between oriented flat charts are rotations and translations.

Let \(SM\) be its unit tangent bundle, \(\nu\) its normalized Liouville
measure, and \(\Psi_t\) its geodesic flow. Folding to the original
triangle sends this flow to the billiard flow. Away from edges, there
are two equally weighted sheets, so the folding map sends \(\nu\) to
\(\mu\). Each regular billiard segment lifts from either initial sheet,
and its time map permutes the two sheet choices.
The set of trajectories hitting a vertex has measure zero.
Indeed, for each fixed finite sequence of crossed sides, unfolding makes
a vertex hit the condition that a straight trajectory pass through a
fixed point. This is a null condition in position--direction space.
The countable union over side sequences, vertices, and both time
directions remains null.

There is no additional obstruction from infinitely many reflections in
finite time. Such an accumulation could occur only at a vertex: an open
side is regular upstairs, and away from the vertices compactness gives
regular neighborhoods. Near a vertex, unfold successive reflections in
the wedge of its fixed positive angle. A straight line avoiding the tip
sweeps at most \(\pi\) in polar angle, so it crosses only finitely many
successive wedge sectors. A line ending at the tip likewise crosses only
finitely many before reaching it. Thus the excluded set accounts for
all finite-time failures of the flow. On its complement the regular
geodesic flow is defined for every time. In flat charts it is translation
at fixed angle and preserves area times angular measure. Its local
preservation, followed along regular trajectories, proves preservation
of \(\nu\); folding gives the same assertion for \(\mu\).
Edge crossings introduce no singularity upstairs.

In a flat chart write \(v=(\cos\theta,\sin\theta)\), and set
\[
 \pa=\tfrac12(\pa_{x_1}-i\pa_{x_2}),\qquad
 \pb=\tfrac12(\pa_{x_1}+i\pa_{x_2}),
\]
\begin{equation}\label{eq:XY}
 X=e^{i\theta}\pa+e^{-i\theta}\pb,\qquad
 Y=i e^{i\theta}\pa-i e^{-i\theta}\pb.
\end{equation}
These are the forward and perpendicular horizontal vector fields on
\(SM\); they commute locally. All distributional equations below are
on the regular bundle. In particular, a bounded function invariant
almost everywhere under each \(\Psi_t\) satisfies \(Xf=0\), by testing
against compactly supported smooth functions.
We use these equalities in \(L^2\) separately for each time, and under
time integrals by Fubini. We do not require a representative invariant
pointwise outside one common exceptional set for all real times.

We use the angular Fourier framework for flat geodesic flows developed
by Forni and Moll \cite[Sections~3.2--3.5]{FM}, deriving the coordinate
rules needed here directly. Fiberwise Fourier projection is globally
defined. Locally write
\[
 P_j u=u_j(x)e^{ij\theta},\qquad j\in\Z.
\]
For a change of chart \(z'=e^{i\phi}z+c\), one has
\(\theta'=\theta+\phi\) and \(u'_j=e^{-ij\phi}u_j\).
Thus the coefficients are sections of flat unitary line bundles.
Their local derivatives transform by
\[
 (\pa u_j)'=e^{-i(j+1)\phi}\pa u_j,\qquad
 (\pb u_j)'=e^{-i(j-1)\phi}\pb u_j.
\]
The flat gradient of a coefficient means its weak gradient in these
trivializations, tested against smooth sections compactly supported
in the regular surface. An \(L^2\) gradient will always mean a finite
global norm on that surface. It does not specify a trace at a cone tip
or membership in a completed operator domain.
We use the normalized area measure for coefficient norms and \(\nu\)
for bundle norms; Parseval then has no extra factor. Flat gradients
are denoted by \(\nabla\). For fields having horizontal derivatives,
\[
 \norm{\nabla u}_2^2=\norm{Xu}_2^2+\norm{Yu}_2^2
       =\sum_j\norm{\nabla u_j}_2^2.
\]
Inner products are complex and linear in the first entry. The mode
identities, also valid distributionally, are
\begin{equation}\label{eq:modes}
 (Xu)_j=\pa u_{j-1}+\pb u_{j+1},\qquad
 (Yu)_j=i\pa u_{j-1}-i\pb u_{j+1}.
\end{equation}
Both terms in each sum have the same unitary transition factor.

\section{Smoothing away from vertices}\label{sec:smoothing}

The next construction allows differentiation and integration by parts
without assigning values at a cone point. Its support margin will be
essential when a smoothed field is extended by zero.

Put \(r(x)=d(x,V)\) in the completed double. There is \(K>0\) such
that, for all sufficiently small \(\eps>0\), every center with
\(r(x)>K\eps\) has an embedded flat disk of radius \(4\eps\).
Near a fixed cone point a disk of radius \(c r(x)\) is embedded for
some \(c>0\); away from the ends, compactness supplies the assertion.
Also, for each fixed \(L\),
\begin{equation}\label{eq:area}
 \operatorname{Area}\{r\le L\eps\}=O_L(\eps^2).
\end{equation}

Choose a nonnegative smooth radial kernel of mass one in the unit disk,
denoted by \(\rho\), and put
\(\rho_\eps(y)=\eps^{-2}\rho(y/\eps)\).
At a center with \(r>K\eps\), define
\[
 (S_\eps u)(x,\theta)
   =\int_{\mathbb R^2}\rho_\eps(y)u(x-y,\theta)\,dy
\]
in an embedded flat disk, keeping the direction parallel. The formula
uses ordinary Euclidean area inside that chart. Radiality makes it
independent of the oriented flat chart. Call \(\{r>K\eps\}\) the
safe region. When a norm is taken over a base set, its inverse image
in \(SM\) is understood.

\begin{lemma}\label{lem:smoothing}
The operators just defined have the following properties.
\begin{enumerate}
\item On the safe region, \(S_\eps\) and \(\eps\nabla S_\eps\) are
uniformly bounded from \(L^2(SM)\) to \(L^2\). They commute locally
with \(X,Y\), including on distributions, and
\(\norm{S_\eps u}_\infty\le\norm{u}_\infty\).
\item For every fixed \(u\in L^2(SM)\),
\begin{equation}\label{eq:small-derivative}
 \eps\norm{Y S_\eps u}_{L^2(r>K\eps)}\longrightarrow0.
\end{equation}
On compact regular sets, \(S_\eps u\to u\) locally in \(L^2\).
\item Fix \(L>K+4\). If \(g\in L^2(SM)\) vanishes on
\(\{r<L\eps\}\), then \(S_\eps g\) is supported in
\(\{r\ge(L-1)\eps\}\) and extends by zero to a field with compact
spatial support in \(M\). For \(u\in L^2(SM)\),
\begin{equation}\label{eq:adjoint}
 \ip{Y S_\eps u}{g}=-\ip{u}{Y S_\eps g}.
\end{equation}
The left pairing uses only safe centers. The analogous identity without
\(Y\) is symmetric. If \(Xg\in L^2\) distributionally, then
\begin{equation}\label{eq:commute}
 X(S_\eps g)=S_\eps(Xg),\qquad
 \norm{X S_\eps g}_2\le C\norm{Xg}_2,
\end{equation}
with the same zero-extension convention.
\end{enumerate}
\end{lemma}

\begin{proof}
In a flat chart containing the convolution disks, the operation is
ordinary spatial convolution at fixed angle. Its kernel, and its first
spatial derivative multiplied by \(\eps\), have uniformly bounded
absolute integrals in either spatial variable. Parallel transport
preserves angular measure. The usual kernel bound on \(L^2\) therefore
proves the asserted norms. Local commutation, positivity, and the
\(L^\infty\) estimate follow from the same chart description.

For a smooth compactly supported input, the expression in
\eqref{eq:small-derivative} tends to zero directly: for small scales its
convolution remains away from the tips and its first derivatives are
uniformly bounded in \(L^2\). Such inputs are
dense in \(L^2(SM)\); the uniform scaled-derivative bound proves the
claim for any fixed input. The ordinary approximate-identity argument
gives local convergence.

For the third assertion, a contributing kernel pair has base points
at distance at most \(\eps\). If one point is in the support of \(g\),
both lie strictly inside the safe region. Reversal of the short segment
reverses its parallel transport and leaves the radial kernel unchanged.
This gives symmetry. In a common flat chart the kernel satisfies
\(Y_x K_\eps(x-y)=-Y_y K_\eps(x-y)\), with the angles identified
by parallel transport. Kernel differentiation and Fubini give
\eqref{eq:adjoint}. No derivative of \(g\) is required.

The support estimate follows from the triangle inequality for \(r\).
In particular, the convolution vanishes in a neighborhood of the
boundary of the safe region, so its extension by zero creates no
boundary derivative. Distributional local commutation proves
\eqref{eq:commute}; \(Xg\) is supported in the same closed spatial
support as \(g\), and the \(L^2\) kernel bound gives the final inequality.
\end{proof}

\section{A localized coefficient estimate}\label{sec:energy}

The estimate below uses the full gradient only where transport invariance
fails. This will confine the cost of cutting off an invariant function
to small cone annuli.
The square summability used in its proof belongs to the regularized
test field; it will not be assumed for the original invariant function.
The equal Cauchy--Riemann norms and parity recurrence also underlie
Forni--Moll's Sobolev arguments \cite[Lemma~3.3, Corollary~4.2,
and Lemma~5.1]{FM}. We prove the support-localized version needed here
on compact subsets of the regular surface.

\begin{lemma}\label{lem:energy}
Let \(w\in L^2(SM)\) have compact spatial support in \(M\), with
\(Xw,Yw\in L^2(SM)\) in the sense of distributions.
If \(Xw\) vanishes over \(M\setminus E\),
where \(E\subset M\) is measurable, then for every \(j\in\Z\),
\begin{equation}\label{eq:energy}
 \norm{\nabla w_j}_2^2
 \le 4\norm{Xw}_2\norm{\ind_E\nabla w}_2.
\end{equation}
\end{lemma}

\begin{proof}
For each coefficient, integration by parts gives
\begin{equation}\label{eq:CR-norms}
 \norm{\pa w_j}_2^2=\norm{\pb w_j}_2^2
                  =\tfrac14\norm{\nabla w_j}_2^2.
\end{equation}
Indeed, the difference of the first two squared norms is proportional
to the integral of \(dw_j\wedge d\overline{w_j}\). The primitive
\(w_j\,d\overline{w_j}\) is global because its constant unitary
transition factors cancel, and its support is compact in the regular
surface. Stokes gives zero. Spatially smooth fields suffice initially;
ordinary local Sobolev approximation gives the stated generality.

Put \(D_j=\norm{\pa w_j}_2^2\). By \eqref{eq:modes} and
\eqref{eq:CR-norms},
\[
 D_j-D_{j+2}
 =\operatorname{Re}\ip{(Xw)_{j+1}}{\pa w_j-\pb w_{j+2}}.
\]
The pairing is over \(E\), since its first factor vanishes elsewhere.
Moreover \(D_k\to0\) as \(k\to+\infty\), because the horizontal
derivatives of \(w\) are square summable in their Fourier modes.
Telescoping in steps of two and applying Cauchy--Schwarz gives
\[
 D_j\le\norm{Xw}_2
 \left(\sum_{k\ge0}
  \norm{\ind_E(\pa w_{j+2k}-\pb w_{j+2k+2})}_2^2\right)^{1/2}.
\]
Pointwise in the base variable, the sum inside is bounded by
\[
 2\sum_{\ell\in\Z}(|\pa w_\ell|^2+|\pb w_\ell|^2)
       =\sum_{\ell\in\Z}|\nabla w_\ell|^2.
\]
Parseval and \eqref{eq:CR-norms} prove \eqref{eq:energy}.
\end{proof}

\begin{lemma}\label{lem:bounded-modes}
On the regular double \(M\) of a fixed nondegenerate triangle,
suppose \(h\in L^\infty(SM)\), \(\norm{h}_\infty\le H\), and
\(Xh=0\) locally in distributions. Then every coefficient has a
global square-integrable flat gradient, and
\begin{equation}\label{eq:bounded-modes}
 \norm{\nabla h_j}_2\le CH\qquad(j\in\Z).
\end{equation}
\end{lemma}

\begin{proof}
Take a smooth spatial cutoff \(\chi_\eps\) equal to zero for
\(r\le(K+1)\eps\), equal to one for \(r\ge(K+2)\eps\), with
\(|\nabla\chi_\eps|\le C/\eps\). For small \(\eps\), its transition
regions are disjoint cone annuli. Set
\(w=\chi_\eps S_\eps h\), extended by zero near the vertices.
This field has compact spatial support and square-integrable horizontal
derivatives. Local commutation gives
\[
 Xw=(X\chi_\eps)S_\eps h.
\]
Thus \(Xw\) is supported in the transition annuli \(E_\eps\),
whose total area is \(O(\eps^2)\). On these annuli the kernel
derivative estimate and boundedness of \(h\) give
\[
 |Xw|+|\nabla w|\le CH/\eps.
\]
Both norms on the right of \eqref{eq:energy} are at most \(CH\).
Hence \(\norm{\nabla w_j}_2\le CH\), independently of \(\eps,j\).

For each fixed \(j\), local \(L^2\) convergence of \(w_j\) to
\(h_j\) and weak compactness of its globally bounded gradients
identify the weak gradient limit with \(\nabla h_j\), by testing
in regular charts. Lower semicontinuity gives
\eqref{eq:bounded-modes}. This argument neither assigns traces at
the cone tips nor asserts square summability over~\(j\).
\end{proof}

The hypothesis of \Cref{lem:bounded-modes} is the local distributional
equation, rather than invariance under the globally defined measurable
flow. This distinction will let us apply it to a weak limit in the next
section without first establishing a new global invariance statement.

\section{From flow invariance to parallel coefficients}\label{sec:rigidity}

We now show that the individual gradients just obtained are zero.
The proof uses the regularized coefficient estimate twice: first to
pass a fixed-mode pairing through the smoothing limit, and then to
pass it through a second approximation limit. The two limits have
different uniform estimates.

\begin{proposition}\label{prop:rigidity}
If \(f\in L^\infty(SM)\) is invariant almost everywhere under
\(\Psi_t\) for each \(t\in\R\), then
\[
 \pa f_j=\pb f_j=0\qquad(j\in\Z)
\]
as distributions on the regular surface.
\end{proposition}

\begin{proof}
Write \(H=\norm{f}_\infty\). The case \(H=0\) is immediate.
We know \(Xf=0\), so \Cref{lem:bounded-modes} applies.
Fix \(m\in\Z\) throughout the following argument.

\medskip\noindent\emph{The endpoint equation.}
Define
\begin{equation}\label{eq:tail}
 U=\sum_{k\ge0}P_{m+2k}f\in L^2(SM),\qquad
 B=e^{i(m-1)\theta}\pb f_m\in L^2(SM).
\end{equation}
The latter expression is a global field of mode \(m-1\).
The recurrence \eqref{eq:modes} implies
\begin{equation}\label{eq:endpoint}
 XU=B
\end{equation}
distributionally. One precise verification uses finite tails \(U_N\):
\[
 XU_N=B+e^{i(m+2N+1)\theta}\pa f_{m+2N}.
\]
The terminal fields \(R_N\) are uniformly bounded in \(L^2\) by
\Cref{lem:bounded-modes}, and their angular modes tend to infinity.
For any fixed \(\varphi\in L^2(SM)\),
\[
 |\ip{R_N}{\varphi}|
 \le CH\norm{P_{m+2N+1}\varphi}_2\longrightarrow0.
\]
Thus they converge weakly to zero; their norms need not tend to zero.
Since
\(U_N\to U\) in \(L^2\), \eqref{eq:endpoint} follows.
On safe disks, smoothing gives
\begin{equation}\label{eq:smoothed-equation}
 X Y S_\eps U=Y S_\eps B.
\end{equation}
Our aim is to prove \(\ip{B}{P_{m-1}Yf}=0\). The recurrence
\eqref{eq:modes} gives \(P_{m-1}Yf=-2iB\), so that identity forces
\(B=0\); the construction below obtains it without assuming
\(Yf\in L^2\).

\medskip\noindent\emph{Long segments with clearance.}
Fix \(s>0\), put \(T=s/\eps\), and fix \(L>K+4\).
This time scale keeps the discarded mass of order \(s\), while the
endpoint error will vanish as \(\eps\to0\) at fixed \(s\).
Let \(G_\eps\) be the set of regular states whose trajectories satisfy
\(r(\Psi_t z)>L\eps\) throughout \(0\le t\le T\).
The distance function is \(1\)-Lipschitz along a unit-speed path.
A time mesh of spacing at most \(\eps\), with at most
\(T/\eps+2\) points, detects every failure of this condition inside
\(\{r\le(L+1)\eps\}\). Measure preservation and \eqref{eq:area}
therefore give
\begin{equation}\label{eq:bad-measure}
 \nu(SM\setminus G_\eps)\le C(T\eps+\eps^2)
                                  =C(s+\eps^2).
\end{equation}
Null sets of singular states may be ignored.

Set \(a_\eps=f\ind_{G_\eps}\) and form the negative-time average
\begin{equation}\label{eq:average}
 g_\eps=\frac1T\int_0^T a_\eps\circ\Psi_{-t}\dd t.
\end{equation}
It has absolute value at most \(H\). If an integrand is nonzero at
\(z\), then \(z\) lies at time \(t\) on a good segment starting
at \(\Psi_{-t}z\). Consequently \(g_\eps=0\) on
\(\{r<L\eps\}\). Invariance of \(f\), Minkowski's inequality,
and \eqref{eq:bad-measure} show that
\begin{equation}\label{eq:average-bounds}
 \norm{g_\eps-f}_2\le CH\sqrt{s+\eps^2},\qquad
 \norm{Xg_\eps}_2\le\frac{2H}{T}.
\end{equation}
For the second assertion, finite-time averaging gives the weak identity
\[
 Xg_\eps=\frac{a_\eps-a_\eps\circ\Psi_{-T}}{T}.
\]
It follows by testing \eqref{eq:average} against a compactly supported
smooth regular field, changing variables by the flow, and integrating
its derivative in \(t\). It does not differentiate the indicator of
\(G_\eps\).

\medskip\noindent\emph{Integrating the smoothed equation.}
We justify use of the fundamental theorem in
\eqref{eq:smoothed-equation} along good segments, despite the absence
of angular smoothing. For each fixed \(\eps\), choose a countable cover
by flat charts enlarged
to contain the required convolution disks. For almost every angular
slice in each chart, the smoothed fields have smooth spatial
representatives and their differential identity holds everywhere in
position. Overlap identities have the same property; include their
exceptional angular sets in this countable collection.

During a visit to a flat chart a regular orbit has constant chart angle.
If it met one of the exceptional slices, it would remain in that slice
for an open time interval, which contains a rational time. For each
chart and rational time the set of such initial states is null, by
measure preservation. Removing their countable union makes the calculus
valid along every chart visit of almost every good segment. Smooth
spatial representatives agree on overlaps, so their identities can be
integrated along the whole segment. This argument also handles its
endpoints by regular continuation there.

Changing variables in the pairing with \eqref{eq:average}, and using
the fundamental theorem, now gives
\begin{equation}\label{eq:segment-identity}
 \ip{Y S_\eps B}{g_\eps}
 =\frac1T\int_{G_\eps}\overline{f(z)}
   \bigl[(Y S_\eps U)(\Psi_T z)-(Y S_\eps U)(z)\bigr]\dd\nu(z).
\end{equation}
Both endpoints are safe. By Cauchy--Schwarz and measure preservation,
the absolute value of the right side is at most
\[
 \frac{2H}{s}\,\eps\norm{Y S_\eps U}_{L^2(r>K\eps)}=o(1)
                 \qquad(\eps\to0,\ s\text{ fixed}),
\]
by \eqref{eq:small-derivative}. Here \(U\) is fixed.

Put \(w_\eps=S_\eps g_\eps\), with the supported extension of
\Cref{lem:smoothing}. Its adjoint identity converts
\eqref{eq:segment-identity} into
\begin{equation}\label{eq:pairing-small}
 \ip{B}{Yw_\eps}=o(1).
\end{equation}
The same lemma and \eqref{eq:average-bounds} give
\[
 \norm{\nabla w_\eps}_2\le CH/\eps,\qquad
 \norm{Xw_\eps}_2\le CH/T.
\]
It has compact spatial support in \(M\). Applying
\Cref{lem:energy} with \(E=M\) yields
\begin{equation}\label{eq:fixed-s-bound}
 \norm{\nabla(w_\eps)_j}_2^2\le CH^2/s\qquad(j\in\Z).
\end{equation}

\medskip\noindent\emph{The first limit: fixed \(s\).}
Take a sequence \(\eps\to0\) and a subsequence along which
\(w_\eps\) converges weakly in \(L^2\) to \(g^{(s)}\).
The closed convex set \(\{|u|\le H\}\subset L^2\) is weakly closed,
so \(\norm{g^{(s)}}_\infty\le H\).
For a fixed smooth compactly supported test field \(\varphi\),
symmetry of convolution gives, for sufficiently small \(\eps\),
\[
 \ip{w_\eps-g_\eps}{\varphi}
       =\ip{g_\eps}{S_\eps\varphi-\varphi}\longrightarrow0.
\]
The uniform \(L^2\) bound and density of these tests show that
\(g_\eps\) has the same weak \(L^2\) limit. Therefore
\begin{equation}\label{eq:limit-function}
 Xg^{(s)}=0,\qquad
 \norm{g^{(s)}-f}_2\le CH\sqrt{s}.
\end{equation}
The first equality follows either from \eqref{eq:average-bounds} or
the bound for \(Xw_\eps\); the second follows by weak lower
semicontinuity.

Only coefficient derivatives at modes \(m\) and \(m-2\) contribute
to \eqref{eq:pairing-small}, because \(B\) has mode \(m-1\).
For fixed \(s\), \eqref{eq:fixed-s-bound} allows a further subsequence
on which both these gradients converge weakly in their global \(L^2\)
spaces. Testing in regular charts identifies the limits as the
gradients of the corresponding coefficients of \(g^{(s)}\).
Since \(B\) is fixed in \(L^2\), passing to the limit gives
\begin{equation}\label{eq:first-limit}
 \ip{B}{P_{m-1}Yg^{(s)}}=0.
\end{equation}
Here this projected derivative means the combination of the two
coefficient derivatives in \eqref{eq:modes}; no full \(L^2\) bound
on \(Yg^{(s)}\) is asserted.

\medskip\noindent\emph{The second limit: \(s\to0\).}
We do not use the deteriorating bound \eqref{eq:fixed-s-bound} for
this limit. Instead, \(g^{(s)}\) is bounded by \(H\) and solves the
exact local equation in \eqref{eq:limit-function}, so
\Cref{lem:bounded-modes} gives
\[
 \norm{\nabla(g^{(s)})_j}_2\le CH,
\]
uniformly in \(s\). Choose \(s_n\downarrow0\) and the limits just
constructed for each \(s_n\). By \eqref{eq:limit-function}, they
converge strongly to \(f\) in \(L^2\). Their relevant gradients
are bounded in \(L^2\), and every weak subsequential limit is
identified locally with the corresponding gradient of \(f\).
Thus \eqref{eq:first-limit} implies
\[
 \ip{B}{P_{m-1}Yf}=0.
\]
The equation \(Xf=0\) in mode \(m-1\) says
\(\pa f_{m-2}=-\pb f_m\). By \eqref{eq:modes},
\[
 P_{m-1}Yf=-2iB,\qquad
 0=\ip{B}{-2iB}=2i\norm{B}_2^2.
\]
Hence \(\pb f_m=0\). Since \(m\) was arbitrary, the recurrence
also gives \(\pa f_j=0\) for every \(j\), proving the proposition.
\end{proof}

\section{Holonomy and the conclusion}\label{sec:conclusion}

\begin{proof}[Proof of \Cref{thm:main}]
Relabel the vertices so that \(\alpha/\pi\) is irrational.
Let \(A\) be a measurable invariant set as in the theorem, and lift
its indicator to a bounded invariant function \(f\) on \(SM\).
By \Cref{prop:rigidity}, each coefficient has zero flat gradient.
In every regular chart it therefore has a constant representative,
and these constants match under the unitary transitions. Each
coefficient is a parallel section.

Continuation around the vertex of angle \(\alpha\) has holonomy
rotation by \(-2\alpha\) modulo \(2\pi\), as follows by developing
its punctured cone of angle \(2\alpha\). A parallel coefficient of
mode \(j\) must consequently be fixed by the phase
\(e^{2ij\alpha}\). If \(j\ne0\), irrationality of \(\alpha/\pi\)
makes that phase different from one, so the coefficient vanishes.
The zeroth coefficient is constant on the connected surface \(M\).
Fourier completeness shows that \(f\) is constant almost everywhere.
As an indicator with the same mass as \(A\), it has integral zero
or one. Thus \(\mu(A)\in\{0,1\}\).
\end{proof}

\begin{remark}
The precise angle condition is visible in the last step. All the analytic
arguments apply to the double of a fixed nondegenerate triangle. A
single irrational angle is enough to eliminate every nonzero parallel
mode, which is why \Cref{thm:main} includes triangles with one rational
angle as well as triangles with three irrational angles. Since the
three angles sum to \(\pi\), exactly one irrational angle and two
rational angles cannot occur. The constants in the analytic estimates
are allowed to depend on the triangle.
\end{remark}

\section{Quantum ergodicity and nodal domains}\label{sec:spectral}
The full-flow theorem has spectral consequences for each fixed triangle
in its scope. The following statement keeps the boundary-condition
dependence of the Cauchy data and their limiting measure explicit.

\begin{corollary}\label{cor:quantum}
Let \(Q\) satisfy the hypotheses of Theorem~\ref{thm:main}, and fix
either Dirichlet (\(B=D\)) or Neumann (\(B=N\)) boundary conditions.
Let \((u_j)_{j\ge1}\) be any complete orthonormal eigenbasis of
\(L^2(Q)\) for the corresponding Laplacian, with
\[
 -\Delta_B u_j=\lambda_j^2u_j,\qquad
 0\le\lambda_1\le\lambda_2\le\cdots,
\]
where frequencies are listed with multiplicity. A set
\(S\subset\mathbb N\) has density one if
\(\#(S\cap\{1,\ldots,N\})/N\to1\).
Then the following conclusions hold.
\begin{enumerate}
\item There is a density-one set \(S_{\mathrm{int}}\) such that, for
every classical order-zero pseudodifferential operator \(A\) whose
Schwartz kernel is compactly supported in \(Q\times Q\),
\[
 \lim_{\substack{j\to\infty\\j\in S_{\mathrm{int}}}}
 \langle Au_j,u_j\rangle_{L^2(Q)}
 =\int_{Q\times S^1}\sigma_0(A)\dd\mu.
\]
Here the Euclidean metric identifies unit covectors with \(Q\times S^1\),
and \(\sigma_0(A)\) is restricted to them.
\item Let \(\Gamma\) be the regular boundary, the union of the three open sides.
For \(\lambda_j>0\), let \(h_j=\lambda_j^{-1}\) and define the Cauchy
data on the open sides by
\[
 u_j^b=\begin{cases}
 h_j\,\partial_\nu u_j|_\Gamma,&B=D,\\
 u_j|_\Gamma,&B=N,
 \end{cases}
\]
where \(\nu\) is the outward unit normal. Write
\(B^*\Gamma=\{(s,\eta)\in T^*\Gamma:|\eta|<1\}\), using arc length
\(s\) and signed tangential covector \(\eta\) on each side. There is
a density-one set \(S_\partial\subset\{j:\lambda_j>0\}\) such that,
for every \(a\in C_c^\infty(T^*\Gamma)\) and a properly supported order-zero
semiclassical quantization \(A_h=\operatorname{Op}_h(a)\),
\[
 \lim_{\substack{j\to\infty\\j\in S_\partial}}
 \langle A_{h_j}u_j^b,u_j^b\rangle_{L^2(\Gamma)}
 =\frac{2}{\pi\operatorname{Area}(Q)}
  \int_{B^*\Gamma}a(s,\eta)w_B(\eta)\dd s\,d\eta,
\]
where
\[
 w_D(\eta)=\sqrt{1-\eta^2},\qquad
 w_N(\eta)=(1-\eta^2)^{-1/2}.
\]
The Cauchy data are not further normalized to have unit boundary norm.
\item If the eigenbasis is real-valued, there is a density-one set
\(S_{\mathrm{nod}}\) such that the number of nodal domains of \(u_j\),
the connected components of \(Q\setminus\{u_j=0\}\), tends to infinity
as \(\lambda_j\to\infty\) with \(j\in S_{\mathrm{nod}}\).
\end{enumerate}
\end{corollary}

\begin{proof}
Place \(\overline Q\) in a sufficiently large flat torus. Its boundary
is Lipschitz, with three smooth open sides and a finite codimension-two
singular set; incident sides meet transversely. Thus the triangle meets
the piecewise-smooth hypotheses in \cite[Section~2]{Burq},
\cite[Definition~3.1]{HZ}, and \cite[Definition~1.16]{Hezari}, including
the last reference's metric-extension condition.
Theorem~\ref{thm:main} supplies the required billiard-flow ergodicity.
The first conclusion is the interior quantum-ergodicity theorem of
Zelditch and Zworski \cite{ZZ}, in the Dirichlet and Neumann formulation
recalled in \cite[Theorem~2 and Remark~2.1]{Burq}.

For the boundary conclusion, discard the null sets of vertex hits and
tangential collisions. The collision section is \(B^*\Gamma\), with
invariant measure \(ds\,d\eta\). Boundary-flight coordinates identify
the flow, up to null sets, with the suspension of its collision map
under the positive free-flight time \(\tau\), and Liouville measure is
proportional to \(ds\,d\eta\,dt\) for \(0\le t<\tau\).
The roof \(\tau\) is positive almost everywhere and bounded by the
diameter of \(Q\). Hence an invariant collision set of positive measure
with positive-measure complement would suspend to a nontrivial invariant
flow set, contradicting Theorem~\ref{thm:main}. The collision map is
therefore ergodic. Hassell--Zelditch's boundary theorem
\cite[Theorem~1.1, equations~(1.3)--(1.5), and Section~7]{HZ} now gives the second
conclusion: its prefactor is
\(4/(\operatorname{vol}(S^1)\operatorname{Area}(Q))
=2/(\pi\operatorname{Area}(Q))\).
Their inward normal instead of our outward normal changes the Dirichlet
data by a sign and leaves the quadratic pairing unchanged.

Finally, Hezari's nodal-domain theorem for two-dimensional ergodic
billiards with piecewise-smooth boundary
\cite[Theorem~1.6]{Hezari} gives the third conclusion for real
eigenfunctions. The finitely many zero-frequency terms have no effect on
density. The three density-one sets may be replaced by their finite
intersection when a simultaneous subsequence is desired.
\end{proof}

The subsequences may depend on the fixed triangle, boundary condition,
and eigenbasis. Zero-density exceptional subsequences remain possible:
these conclusions assert neither quantum unique ergodicity nor an
equidistribution or nodal-count rate, mixing, or a spectral-statistics
law.


\begingroup
\raggedright
\bibliographystyle{plain}
\bibliography{references}
\endgroup
\end{document}
